\documentclass[10pt,a4paper]{amsart}
\usepackage[margin=19mm]{geometry}
\usepackage{amsmath,amssymb,mathtools}
\usepackage[scr=rsfs,cal=euler]{mathalfa}
\usepackage{xfrac}
\usepackage[unicode,hidelinks,bookmarks=false]{hyperref}
\newcommand{\ie}{i.e.\ }
\newcommand{\cf}{cf.\ }
\newcommand{\resp}{respectively\ }
\newcommand{\Subsection}{\subsection}
\providecommand{\nobreakdash}{\nobreak\hbox{-}}
\setlength{\emergencystretch}{2em}
\multlinegap0pt
\usepackage{amssymb}
\usepackage{mathtools}
\usepackage{bm}
\usepackage[shortlabels]{enumitem}
\newtheorem{corollary}{Corollary}
\newtheorem{theorem}{Theorem}
\newtheorem{lemma}{Lemma}
\title{Estimation of Riemannian Quantities from Noisy Data via Density Derivatives}
\author{Junhao Chen, Ruowei Li, Zhigang Yao}
\date{Source: arXiv:2603.27171v1 (CC BY 4.0)}
\begin{document}
\maketitle

\noindent\emph{Source and licence.} Estimation of Riemannian Quantities from Noisy Data via Density Derivatives. Version arXiv:2603.27171v1 (CC BY 4.0), distributed by arXiv under the Creative Commons Attribution 4.0 International licence, http://creativecommons.org/licenses/by/4.0/, as recorded in the arXiv licence metadata for this exact version and shown on the abstract page https://arxiv.org/abs/2603.27171v1. Full licence notice: \texttt{licenses/arxiv-2603.27171-CC-BY-4.0.txt}.\par\smallskip

\noindent\textit{Editorial scope.} Counted unit: the article's theorem th:tangent-space (source0.tex lines 389--395). The article's complete original proof of it is delivered with it (source0.tex lines 1936--2139, one \texttt{proof} environment of 204 lines, which the article places away from the statement). No mathematics is written, repaired or re-derived: the statement and the proof are the article's own lines, in the article's own notation, and no retained line is substituted or re-flowed.  Retained uncounted as the source-local context the proof needs: lines 288--311; lines 346--355; lines 358--365; lines 464--469.  Nothing was omitted from any retained span, so the package carries no omission record.  No retained line contains a citation, so the wrapper carries no bibliography and no bibliography entry.  No retained line contains a figure environment, an image-inclusion command or drawing code, so no image is delivered and the package declares no asset dependency.  Editorial formatting only, in addition: an AMS \texttt{amsart} preamble that restates the article's own declarations and macros used by the retained lines, namely the environments \texttt{corollary}, \texttt{theorem}, \texttt{lemma} on separate counters, together with the standard packages those lines need.  The retained environments print in this document as Corollary 1, Theorem 1, Lemma 1 in order of appearance, whereas the article prints the same items with the numbering of its own full document.  The complete native source of the article as distributed in the arXiv e-print for this exact version is bundled unchanged as \texttt{sources/197328/source0.tex}.
\begin{corollary}\label{co:log-density}
Uniformly for $y\in T(\tau-\varepsilon)$ and $0<\sigma\le\sigma_0$,
\begin{equation}\label{eq:log density}
\log P_\sigma(y)
=
\log \big(V_M(2\pi\sigma^2)^{\frac{D-d}{2}}\big)^{-1}
-
\frac{\|v_y\|^2}{2\sigma^2}
-
\frac{1}{2}\log\det A_y
+
R(y,\sigma),
\end{equation}
where $R(y,\sigma)$ is $C^3$ on $T(\tau-\varepsilon)$ and its ambient derivatives satisfy
\[
|R(y,\sigma)| \le C_{0}\big(\sigma\|v_y\|+\sigma^2\big),
\quad
\|\nabla_y^m R(y,\sigma)\|_{\mathrm{op}}
\le
C_{m}\sigma,
\qquad
m\in\{1,2,3\}.
\]
\end{corollary}

\begin{theorem}\label{th:noisy-hessian}
Uniformly for $y\in T(\tau-\varepsilon)$ and $0<\sigma\le\sigma_0$,
\[
\mathcal{H}_\sigma(y)
= - \frac{1}{\sigma^2}P_N(y)
- \frac{1}{\sigma^2}(I_{T_{\pi(y)}M} - A_y^{-1})P_T(y)
+ O(1),
\]
where the $O(1)$ term is taken in the operator norm.
\end{theorem}

\begin{corollary}\label{co:noisy-hessian-gap}
For each fixed $\varepsilon\in(0,\tau)$, there exist constants $c_\varepsilon>0$ and $\sigma_0>0$ such that the following holds. Let $\lambda_1(y,\sigma)\ge \cdots \ge \lambda_D(y,\sigma)$
be the eigenvalues of $\mathcal H_\sigma(y)$. Then, for every $y\in T(\tau-\varepsilon)$ and every $0<\sigma\le \sigma_0$,
\[
\lambda_d(y,\sigma)-\lambda_{d+1}(y,\sigma)\ge c_\varepsilon \sigma^{-2}.
\]
The same spectral gap holds for $\mathcal H_0(y)$.
\end{corollary}

\begin{lemma}\label{lem:second-fundamental-form-general}
    Let $P_T(y)$ be a $C^1$ tangential projection field on $\mathcal T(\tau)$. Then, for every $x\in M$ and $u,v\in T_xM$,
    \begin{align}
        \Pi_x(u,v) = \nabla_uP_T(x) v.
    \end{align}
\end{lemma}

\begin{theorem}\label{th:tangent-space}
Let $\widehat T_yM$ be the eigenspace associated with the $d$ largest eigenvalues of $\mathcal H_\sigma(y)$. Then there exist constants $C$ and $\sigma_0>0$ such that, for every $y\in T(\tau-\varepsilon)$ and every $0<\sigma\le\sigma_0$,
\[
\sin\{\Theta(\widehat{T}_{y}M,T_{\pi(y)}M)\} \le C\sigma^2.
\]
Moreover, for any sequence with $\|v_y\|\to0$ and $\sigma\to0$, the largest consecutive spectral gap of $\mathcal H_\sigma(y)$ is attained at $k=d$. Consequently, $\widehat d=d$ along the sequence.
\end{theorem}

\begin{proof}
Write $x:=\pi(y)$. Recall that $P_T(y)$ denotes the orthogonal projection onto $T_xM$, while $\widehat P_T(y)$ denotes the spectral projector onto the $d$ largest eigenvalues of $\mathcal{H}_\sigma(y)$. By Theorem~\ref{th:tangent-space},
\[
\|\widehat P_T(y)-P_T(y)\|_{\mathrm{op}}
\le C\sigma^2.
\] 
By Lemma~\ref{lem:second-fundamental-form-general} applied at $x$ we have,
\[
\Pi_x(u,v) = \nabla_uP_T(x)v.
\]
The estimator in the theorem is
\[
\widehat\Pi_y(u,v)
=
\nabla_u\widehat P_T(y)v.
\]
We decompose
\[
\widehat\Pi_y(u,v)-\Pi_x(u,v)
=
\big(\nabla_u\widehat P_T(y)-\nabla_uP_T(y)\big)v
+
\big(\nabla_uP_T(y)-\nabla_uP_T(x)\big)v.
\]
%Lip bound
For the second term, since $P_T(y)$ is $C^{k-1}$ on $T(\tau-\varepsilon)$, $\nabla P_T(y)$ is Lipschitz on $T(\tau-\varepsilon)$, i.e.,
\[
\|\nabla_uP_T(y)-\nabla_uP_T(x)\|_{\mathrm{op}}
\le C\|v_y\|.
\]
%asymptotic error
Next we consider the first term. Let
\[
P_N(y):=I_D-P_T(y),
\qquad
\mathcal H_0(y):= -\sigma^{-2}P_N(y)-\sigma^{-2}(I_{T_xM}-A_y^{-1})P_T(y).
\]
Then $P_T(y)$ is exactly the spectral projector of $\mathcal H_0(y)$ associated with the largest $d$ eigenvalues.

Let
\[
\widetilde {\mathcal H}_0(y):=\sigma^2 \mathcal H_0(y),
\qquad
\widetilde {\mathcal H}_\sigma(y):=\sigma^2 \mathcal H_\sigma(y).
\]
By the uniform boundedness of $A_y^{-1}$ on $T(\tau-\varepsilon)$ and the expansion of $\mathcal H_\sigma$ in Theorem~\ref{th:noisy-hessian}, the spectra of $\widetilde {\mathcal H}_0(y)$ and $\widetilde {\mathcal H}_\sigma(y)$ remain in a compact interval independent of $\sigma$. Moreover, by Corollary~\ref{co:noisy-hessian-gap}, their tangential and normal spectral clusters are separated by a uniform positive gap.
Hence one may choose a fixed positively oriented contour $\widetilde\Gamma$ that encloses the tangential cluster and excludes the normal cluster for both $\widetilde {\mathcal H}_0(y)$ and $\widetilde {\mathcal H}_\sigma(y)$.
Setting
\[
\Gamma:=\sigma^{-2}\widetilde\Gamma,
\]
we obtain
\[
\operatorname{length}(\Gamma)\le C_1\sigma^{-2},
\qquad
\operatorname{dist}\big(\Gamma,\operatorname{spec}(\mathcal H_0(y))\big)\ge C_2\sigma^{-2},
\qquad
\operatorname{dist}\big(\Gamma,\operatorname{spec}(\mathcal H_\sigma(y))\big)\ge C_3\sigma^{-2}.
\]
Therefore
\[
\sup_{z\in\Gamma}\|(\mathcal H_0(y)-zI)^{-1}\|_{\mathrm{op}}
+
\sup_{z\in\Gamma}\|(\mathcal H_\sigma(y)-zI)^{-1}\|_{\mathrm{op}}
\le C\sigma^2.
\]
For $\mathcal H$ with spectrum separated by $\Gamma$, define the corresponding Riesz projector by
\[
P[\mathcal H]=-\frac{1}{2\pi i}\int_{\Gamma}(\mathcal H-zI)^{-1}dz.
\]
Its derivative in the direction $K$ is
\[
DP[\mathcal H](K)=\frac{1}{2\pi i}\int_{\Gamma}(\mathcal H-zI)^{-1}K(\mathcal H-zI)^{-1}dz.
\]
Therefore,
\[
\nabla_u\widehat P_T(y)
=
DP[\mathcal H_\sigma(y)]\big(\nabla_u\mathcal H_\sigma(y)\big),
\qquad
\nabla_uP_T(y)
=
DP[\mathcal H_0(y)]\big(\nabla_u\mathcal H_0(y)\big).
\]
Set
\[
\Delta \mathcal H:=\mathcal H_\sigma(y)-\mathcal H_0(y),
\qquad
\Delta K:=\nabla_u\mathcal H_\sigma(y)-\nabla_u\mathcal H_0(y),
\qquad
K_0:=\nabla_u\mathcal H_0(y).
\]
Then
\[
\nabla_u\widehat P_T(y)-\nabla_uP_T(y)
=
DP[\mathcal H_\sigma(y)](\Delta K)
+
\big(DP[\mathcal H_\sigma(y)]-DP[\mathcal H_0(y)]\big)(K_0).
\]
Let $R_\sigma(z)=(\mathcal H_\sigma(y)-zI)^{-1}$ and
$R_0(z)=(\mathcal H_0(y)-zI)^{-1}$. By the above choice of $\Gamma$,
\[
\sup_{z\in\Gamma}\|R_\sigma(z)\|_{\mathrm{op}}
,\ 
\sup_{z\in\Gamma}\|R_0(z)\|_{\mathrm{op}}
\le C\sigma^2.
\]
Hence
\[
\|DP[\mathcal H_\sigma(y)](\Delta K)\|_{\mathrm{op}}
\le
C\operatorname{length}(\Gamma)
\sup_{z\in\Gamma}\|R_\sigma(z)\|_{\mathrm{op}}^2
\|\Delta K\|_{\mathrm{op}}
\le
C\sigma^2\|\Delta K\|_{\mathrm{op}}.
\]
Moreover, by the resolvent identity,
\[
R_\sigma(z)-R_0(z)
=
-R_\sigma(z)\Delta \mathcal HR_0(z),
\]
so that
\[
\sup_{z\in\Gamma}\|R_\sigma(z)-R_0(z)\|_{\mathrm{op}}
\le
C\sigma^4\|\Delta \mathcal H\|_{\mathrm{op}}.
\]
Therefore,
\begin{align*}
&\big\|
\big(DP[\mathcal H_\sigma(y)]-DP[\mathcal H_0(y)]\big)(K_0)
\big\|_{\mathrm{op}}\\
&\le
C\sup_{z\in\Gamma}(\|R_\sigma(z)\|+\|R_0(z)\|)\sup_{z\in\Gamma}\|R_\sigma(z)-R_0(z)\|_{\mathrm{op}}\|K_0\|_{\mathrm{op}}
\le
C\sigma^4\|\Delta \mathcal H\|_{\mathrm{op}}\|K_0\|_{\mathrm{op}}.
\end{align*}
We now estimate the three quantities above. First, by Theorem~\ref{th:noisy-hessian},
\[
\|\Delta \mathcal H\|_{\mathrm{op}}
=
\|\mathcal H_\sigma(y)-\mathcal H_0(y)\|_{\mathrm{op}}
=
O(1).
\]
Next,
\[
K_0
=
\nabla_u\mathcal H_0(y)
=
-\sigma^{-2}\Big(\nabla_uP_N(y)+\nabla_u \big((I_{T_xM}-A_y^{-1})P_T(y)\big)\Big),
\]
and since $\nabla P_T$, $\nabla P_N$, and $\nabla A_y^{-1}$ are uniformly bounded on $T(\tau-\varepsilon)$,
\[
\|K_0\|_{\mathrm{op}}\le C\sigma^{-2}.
\]
Finally, we claim that $\|\Delta K\|_{\mathrm{op}} \le C$. From the expansion in the proof of
Theorem~\ref{th:noisy-hessian},
\[
\mathcal H_\sigma(y)
=
-\sigma^{-2}P_N(y)
-\sigma^{-2}(I_{T_xM}-A_y^{-1})P_T(y)
-\frac12\nabla^2\log\det A_y
+\nabla^2R(y,\sigma).
\]
Thus
\[
\Delta \mathcal H 
=
-\frac12\nabla^2\log\det A_y
+\nabla^2R(y,\sigma).
\]
Differentiating in the tangential direction $u\in T_xM$, the above terms satisfy
\[
\|\nabla_u\nabla^2\log\det A_y\|_{\mathrm{op}} \le C,
\qquad
\|\nabla_u\nabla^2R(y,\sigma)\|_{\mathrm{op}} \le C\sigma,
\]
uniformly on $T(\tau-\varepsilon)$, by smoothness of $A_y$ and
Corollary~\ref{co:log-density}. Therefore,
\[
\|\Delta K\|_{\mathrm{op}}
=
\|\nabla_u\mathcal H_\sigma(y)-\nabla_u\mathcal H_0(y)\|_{\mathrm{op}}
\le C.
\]
Substituting the bounds for $\Delta \mathcal H$, $\Delta K$, and $K_0$ into the previous estimates, we obtain
\[
\|\nabla_u\widehat P_T(y)-\nabla_uP_T(y)\|_{\mathrm{op}}
\le C\sigma^2.
\]
Combining this with the estimate for the second term yields
\[
\|\widehat\Pi_y-\Pi_x\|_{\mathrm{op}}
\le
C(\|v_y\|+\sigma^2).
\]
Since $x=\pi(y)$, this proves the claim.
\end{proof}

\end{document}
